
\begin{frame}

\begin{center}
\large
\textbf{
An Automatic Finite-Sample Robustness Metric:
\\Can Dropping a Little Data Make a Big Difference?}
\end{center}

\hrulefill
\vspace{1em}

Ryan Giordano (\texttt{rgiordano@berkeley.edu})\footnotemark[1]\\[1em]

JSM 2026:
Data Selection Beyond Estimation Efficiency
%: Robustness and Energy-Efficient Model Training


\footnotetext[1]{With coauthors Rafe Meager (UNSW) and Tamara Broderick (MIT)}

\end{frame}

%%%%%%%%%%%%%%%%%%%%%%%%%%%%%%%%%%%%%%%%%%%%%%%%%%%%%%%%%%%%%%%%%%%%%%%%
%%%%%%%%%%%%%%%%%%%%%%%%%%%%%%%%%%%%%%%%%%%%%%%%%%%%%%%%%%%%%%%%%%%%%%%%
%%%%%%%%%%%%%%%%%%%%%%%%%%%%%%%%%%%%%%%%%%%%%%%%%%%%%%%%%%%%%%%%%%%%%%%%

\begin{frame}[t]{Dropping data: Mexico Microcredit}

% The success of statistics has led to data from particular populations
% being used increasingly to draw general conclusions.
% %
% Drawing the right conclusions affects peoples' lives!

\vspace{1em} \textbf{Example:} \citet{angelucci2015microcredit}, a randomized
controlled trial study of the efficacy of microcredit in Mexico based on $N =
$16,560 data points. A regression was run to estimate the average effect of
microcredit.

\hrulefill

\vspace{1em}
\textbf{Original result: }
Treatment effect statistically insignificant at 95\%.

\vspace{1em}
\textbf{Policy implication: } Disinvest in microcredit initiatives.

\hrulefill

\pause

\vspace{1em} \textbf{Data dropping: } Can produce both positive and negative
statististically significant results dropping no more than 15 data points
($<$ 0.1\%).

\vspace{1em}
\textbf{Policy implication: } Run a higher-powered study (not just
larger $N$).


\hrulefill

\pause

Cannot find influential subsets by brute force!

\vspace{1em}
\textbf{We provide a fast, automatic tool to  approximately identify the
most influential set of points.}


\end{frame}


% %%%%%%%%%%%%%%%%%%%%%%%%%%%%%%%%%%%%%%%%%%%%%%%%%%%%%%%%%%%%%%%%%%%%%%%%
% %%%%%%%%%%%%%%%%%%%%%%%%%%%%%%%%%%%%%%%%%%%%%%%%%%%%%%%%%%%%%%%%%%%%%%%%
% %%%%%%%%%%%%%%%%%%%%%%%%%%%%%%%%%%%%%%%%%%%%%%%%%%%%%%%%%%%%%%%%%%%%%%%%


\begin{frame}{Outline}
\begin{itemize}
    \item How do we identify influential sets?  When is our method accurate?
    \item Data dropping as distributional robustness
    \item Do we identify ``the most important datapoints?''
\end{itemize}
\end{frame}

% % %%%%%%%%%%%%%%%%%%%%%%%%%%%%%%%%%%%%%%%%%%%%%%%%%%%%%%%%%%%%%%%%%%%%%%%%
% % %%%%%%%%%%%%%%%%%%%%%%%%%%%%%%%%%%%%%%%%%%%%%%%%%%%%%%%%%%%%%%%%%%%%%%%%
% % %%%%%%%%%%%%%%%%%%%%%%%%%%%%%%%%%%%%%%%%%%%%%%%%%%%%%%%%%%%%%%%%%%%%%%%%

% \begin{frame}{Dropping data: Motivation}

% When do you care that you can \textbf{reverse your conclusion} by removing a
% \textbf{small proportion} of your data?

% \pause \vspace{1em} Not always!  But sometimes you surely do, especially when
% you want to \textbf{generalize to unseen, systematically different populations}.

% \vspace{1em}
% Suppose you have a farm, and want to know whether
% your average yield is $>$ 170 bushels per acre.
% At harvest, you measure 200 bushels per acre.

% \pause
% \begin{itemize}
%     \item Scenario one: $>$ 170 bushels
%         per acre means you make a profit.
%         \begin{itemize}
%             \item Don’t care about sensitivity to small subsets.
%         \end{itemize}
%     \pause
%     \item Scenario two: Want to recommend methods to a distant friend.
%     \begin{itemize}
%         \item Might care about sensitivity to small subsets!
%     \end{itemize}
% \end{itemize}

% \pause

% \vspace{1em}
% Specifically, often in statistical applications:
% %
% \begin{itemize}
% \item Policy population is different from analyzed population,
% \item Small fractions of data are missing not-at-random,
% \item We report a convenient summary (e.g. mean) of a complex effect.
% % \item Models are stylized proxies of reality.
% \end{itemize}


% \end{frame}
